\section{Limitations}\label{sec:limits}
% Objective: Surface two days late; realised not implied correlation; Treasury not OIS; barrier coverage; local volatility's forward skew; six notes.

\subsection{Inputs that are not what a desk would use}

The volatility surface is observed two trading days after five of the six notes
priced, which is a consequence of the filing window in \S\ref{sec:data} rather
than a choice. \S\ref{sec:inputs} measures the cost of that gap, note by note,
and it runs the wrong way to help: correcting for it lowers values further.

Correlation is realised rather than implied, for the reason given in
\S\ref{sec:method}: vanilla options carry no co-movement information and no free
source of implied correlation exists. The bias runs in a known direction, and
\S\ref{sec:inputs} treats it as one of the surviving candidates rather than
burying it here.

Discounting uses the Treasury par yield curve where a dealer would use SOFR
OIS, again because no free source of OIS quotes exists. The two differ by a
handful of basis points at these tenors. That is small against funding spreads
measured in the hundreds, and it is not nothing.

\subsection{The Dow leg, and what the coverage test does not promise}

One leg is priced off fund options standing in for an index
(\S\ref{sec:data}). Two problems follow, and the second is worse than the
first. The level bias is measured at 0.52 volatility points and is worth \$2.61
on the affected note. The more serious problem is that the fund's long-dated
quotes never reach the strike where that note's barrier sits, so for the whole
of its life past the first year the downside is priced off an extrapolation
rather than a quote. That note is reported throughout and excluded from every
quoted range, and its \$0.37 standard error understates its true uncertainty by
an amount this study cannot quantify.

The coverage test that produced that exclusion is a threshold, not a guarantee.
It asks whether a quoted strike exists at or below the trigger on fitted slices
past one year, and a thin market can satisfy it on a single strike while still
pricing the barrier off very little. The other five notes reference index
options deep enough that this is unlikely. Nothing here rules it out.
Appendix~\ref{app:coverage} gives the counts leg by leg so a reader can apply a
stricter threshold.

\subsection{What local volatility does not do}

The engine is calibrated to European option prices and reproduces them, within
the bounds \S\ref{sec:validation} reports. These notes are path-dependent, and
nothing in that validation establishes that local volatility prices the path
dependence correctly. It establishes that the terminal distributions are
right.

This is the model's known weakness rather than an incidental one. Local
volatility reproduces today's surface exactly and misgoverns how that surface
evolves, flattening forward skew in a way that stochastic and stochastic-local
volatility models are built to repair \citep{Qu2016}. It remains market
practice for exotics of this shape because it handles multi-dimensionality,
where stochastic volatility models become expensive to calibrate across several
correlated underlyings \citep{Grzelak2019}. For a worst-of autocallable, whose
value depends on the joint distribution at every observation date and not only
at maturity, the dynamics are a real exposure.

The flat-volatility comparison in \S\ref{sec:inputs} bounds how much the
\emph{shape} of the surface matters. It does not bound how much the dynamics of
that shape matter, and no result in this paper does. What can be said is that
the direction is not unknown: \citet{Farkas2022} report that a Heston
local-stochastic volatility model prices autocallables above local volatility,
with the stochastic parameters moving the autocallable price while the fit to
European options is unchanged. Their products are single-asset. Whether the
same holds for a worst-of on three underlyings, where \citet{DeCol2017} find
the correlations between stochastic variances to be central, is not something
this paper establishes.

\subsection{Exercise, precision and sample}

The issuer's call is solved under optimal exercise, and Longstaff--Schwartz
with a finite basis returns a policy that is suboptimal to some degree. That
error biases the investor's value upward relative to true optimality, so it
pushes the model value toward the disclosed figure and cannot account for a
model that is already too low. The direction is favourable to the finding,
which is worth stating explicitly rather than leaving as an unexamined
approximation.

One issuer discloses its estimated value to the nearest dollar rather than the
cent, so the comparison for that note carries a dollar of irreducible
disclosure precision against a gap of \$19.16. It does not affect the sign of
anything reported.

Six notes, four issuers, one pricing week. Every number here is a case study
and none is an estimate of a population parameter. The 129 final supplements
identified for that week would support a cross-sectional analysis of the
disclosed margin against tenor, coupon, barrier distance and the number of
underlyings. That is left undone, and it is the obvious next thing to do.
